% ID: SEL-D0069-OA-E01
% ESTADO: cotejada. Segunda pasada visual y borrador compilado revisados.
% Fuente: Selectividad/sel_2009_matematicas_aplicadas/mat apli CC SS03.pdf
% SHA-256: 8cf5f487f300c573430722e40f741601da1756c5e8c13e541f697d1620373545
% Página: 1; opción A; ejercicio original 1.

Sea la igualdad $A\cdot X+B=A$, donde $A$, $X$ y $B$ son matrices cuadradas de la misma dimensión.
\begin{enumerate}[label=\alph*)]
\item \textbf{(1 punto)} Despeje la matriz $X$ en la igualdad anterior, sabiendo que $A$ tiene inversa.
\item \textbf{(2 puntos)} Obtenga la matriz $X$ en la igualdad anterior, siendo
\[
A=\begin{pmatrix}2&5\\1&3\end{pmatrix}\ \text{y}\ B=\begin{pmatrix}0&-3\\-1&2\end{pmatrix}.
\]
\end{enumerate}
